\section{Conclusiones}

A partir del an\'alisis independiente de transferencia de calor del tanque de almacenamiento de glucosa Tag 53A-90A-0056, se establecen las siguientes conclusiones t\'ecnicas.

\subsection{Transferencia de Calor}

Los hallazgos principales del an\'alisis t\'ermico demuestran que el sistema de calentamiento est\'a gobernado por la resistencia convectiva del producto. La Tabla~\ref{tab:conclusiones_termicas} resume las conclusiones fundamentales de esta secci\'on.

\begin{table}[H]
\centering
\caption{Conclusiones del an\'alisis de transferencia de calor.}
\label{tab:conclusiones_termicas}
\setcounter{itemcount}{0}
\begin{tabularx}{\textwidth}{@{}lX@{}}
\toprule
\textbf{Par\'ametro} & \textbf{Conclusi\'on T\'ecnica} \\
\midrule
Mecanismo Limitante & La convecci\'on natural de la glucosa representa entre el 97,7~\% y el 98,8~\% de la resistencia t\'ermica total debido a su elevada viscosidad ($\sim$212{.}000~cP a 25~\textdegree C). \\
Eficacia del Sistema & El coeficiente global $U$ se sit\'ua en el rango de 21 a 36~W/m\textsuperscript{2}$\cdot$\textdegree C. El n\'umero de Prandtl ($\mathrm{Pr} \sim 10^{4}$) refleja la dominancia de la difusi\'on viscosa sobre la t\'ermica, estableciendo la convecci\'on natural como mecanismo limitante del proceso, con independencia de la geometr\'ia de la chaqueta. \\
Optimizaci\'on de Gradiente & El uso de agua a 75~\textdegree C reduce el tiempo de preparaci\'on log\'istica en un 46~\% frente a la operaci\'on a 65~\textdegree C, permitiendo alcanzar los 57~\textdegree C operativos en 415,2~h (Escenario~3). \\
Eficiencia Energ\'etica & El aislamiento de lana mineral de 2~pulgadas reduce las p\'erdidas al ambiente en un 91,6~\%, garantizando una temperatura superficial de seguridad de 28,3~\textdegree C y un balance neto positivo de 15,3~MJ/h para el flujo medio de 5.000~kg/h a 60~\textdegree C. \\
Validaci\'on CFD & La simulaci\'on en COMSOL Multiphysics~6.4 valid\'o el modelo anal\'itico con una concordancia del 95~\% ($U_{\text{CFD}} = 38$~W/m\textsuperscript{2}$\cdot$\textdegree C vs.\ $U_{\text{anal\'itico}} = 36,2$~W/m\textsuperscript{2}$\cdot$\textdegree C), confirmando la precisi\'on de las estimaciones transitorias. \\
\bottomrule
\end{tabularx}
\end{table}

Los resultados primarios referentes a los tiempos de calentamiento requeridos para la disponibilidad operativa del fluido se consolidan en la Tabla~\ref{tab:resumen_tiempos}.

\begin{table}[H]
\centering
\caption{Consolidado de resultados operacionales por escenario.}
\label{tab:resumen_tiempos}
\setcounter{itemcount}{0}
\begin{tabularx}{\textwidth}{@{}lccccX@{}}
\toprule
\textbf{Escenario} & \textbf{Llenado} & \textbf{$T_{agua}$ [\textdegree C]} & \textbf{$T_{fin}$ [\textdegree C]} & \textbf{Tiempo Total} & \textbf{Comportamiento de Descarga} \\
\midrule
\textbf{Escenario 1} & 24~m\textsuperscript{3} (7,6~\%) & 65 & 60 & $\sim$110~h & Incremento ligero durante descarga de 2~h. \\
\textbf{Escenario 2} & 178~m\textsuperscript{3} (80~\%) & 65 & 57 & 772,8~h & Arranque desde fr\'io con agua a 65~\textdegree C. \\
\textbf{Escenario 3} & 178~m\textsuperscript{3} (80~\%) & 75 & 57 & 415,2~h & Reducci\'on del 46~\% vs.\ Escenario~2; cota m\'axima recomendada a 472,8~h. \\
\textbf{Escenario 4} & Capacidad diaria 120~ton & 75 & 57 & Ciclo diario & 5 descargas/d\'ia; flujo medio 5,0~ton/h dentro de capacidad de mantenimiento. \\
\textbf{Escenario 5} & Inventario caliente & 75 & 57 & Ciclo diario & 5 descargas con recalentamiento entre descargas; temperatura estable en 56,6--57,0~\textdegree C. \\
\bottomrule
\end{tabularx}
\end{table}

\subsubsection{Comparativa de Configuraciones de Chaqueta}

El an\'alisis comparativo entre la chaqueta dimple actual ($A_{actual} = 28$~m\textsuperscript{2}) y la media ca\~na del fondo construido ($A_{reemp} = 13$~m\textsuperscript{2}) establece con precisi\'on cuantitativa el impacto de la reducci\'on de \'area sobre la capacidad de calentamiento del sistema. La Tabla~\ref{tab:conclusiones_chaquetas} consolida los hallazgos principales.

\begin{table}[H]
\centering
\caption{Conclusiones de la comparativa de configuraciones de chaqueta ($T_w = 75$~\textdegree C, $v = 2,5$~m/s, tanque al 80~\%).}
\label{tab:conclusiones_chaquetas}
\setcounter{itemcount}{0}
\begin{tabularx}{\textwidth}{@{}lX@{}}
\toprule
\textbf{Par\'ametro} & \textbf{Conclusi\'on} \\
\midrule
Invarianza de $U$ & El coeficiente global $U$ es id\'entico para ambas configuraciones ($\Delta U / U < 0,3$~\%); la resistencia del lado de la glucosa concentra el 97,7--98,8~\% de $R_{total}$, con independencia de la geometr\'ia de la chaqueta. \\
Ratio de Capacidad \textit{Batch} & $t_{reemp}/t_{actual} = 2,15 \approx A_{actual}/A_{reemp} = 2,15$. Calentamiento 25--57~\textdegree C: actual 170,7~h vs.\ fondo de reemplazo 367,5~h. Calentamiento 55--57~\textdegree C: actual 15,5~h vs.\ fondo de reemplazo 33,2~h. \\
Flujo Continuo desde 55~\textdegree C & Chaqueta actual: flujo m\'aximo $>$ 16~ton/h. Chaqueta del fondo de reemplazo: flujo m\'aximo 7,54~ton/h. Para operaci\'on de mantenimiento a flujos superiores a 7,5~ton/h, la chaqueta del fondo de reemplazo resulta insuficiente. \\
Flujo Continuo desde 25~\textdegree C & Ambas configuraciones son insuficientes para flujos industriales: la chaqueta actual alcanza solo 1,47~ton/h y la del fondo de reemplazo solo 0,69~ton/h. El calentamiento desde fr\'io en modo continuo requiere agitaci\'on mec\'anica para ser viable a escala industrial. \\
Palanca de Optimizaci\'on & La selecci\'on del tipo de chaqueta (dimple vs.\ media ca\~na) es irrelevante para el coeficiente $U$. La \'unica palanca en el lado de la chaqueta es el \'area. La mejora sustancial requiere actuar sobre $h_o$ mediante agitaci\'on mec\'anica. \\
\bottomrule
\end{tabularx}
\end{table}

En s\'intesis, el sistema de calentamiento por media ca\~na de 13~m\textsuperscript{2} es t\'ecnicamente capaz de sostener el ciclo operativo oficial de cinco descargas diarias de 24~toneladas, siempre que la glucosa de alimentaci\'on ingrese al tanque dentro del rango 57--60~\textdegree C, se opere con agua de calentamiento a 75~\textdegree C y se mantenga instalado el aislamiento t\'ermico de 50,8~mm. El arranque desde fr\'io del tanque completo no es viable en un ciclo diario, por lo que la operaci\'on debe basarse en el mantenimiento de un inventario caliente.
